\begin{table}
\caption{Tabela 1. Portão dos descritores nas séries de 8 dias (1,104 pontos, 20% de ruído de observação, 20 repetições): fração de rejeição das substitutas IAAFT nos controles caóticos (poder) e nos aleatórios (tamanho). Família primária: os descritores que entram na leitura, fixados no pré-registro das séries de vegetação.}
\label{tab:12_floresta_atmosfera::portao}
\begin{tabular}{llllr}
\toprule
Descritor & Família & Poder (logística, Hénon, Lorenz) & Rejeição (branco, AR(1)) & Rejeição (sazonal com tendência) \\
\midrule
Assimetria temporal & primária & 1,00, 1,00, 0,00 & 0,10, 0,05 & 0,1 \\
Cao (E2) & primária & 1,00, 1,00, 1,00 & 0,10, 0,05 & 0,1 \\
Dimensão de correlação & exploratória & 0,15, 0,30, 0,10 & 0,15, 0,00 & 0,1 \\
Lyapunov (Kantz) & exploratória & 1,00, 1,00, 1,00 & 0,05, 0,05 & 1 \\
Lyapunov (Rosenstein) & exploratória & 1,00, 0,95, 1,00 & 0,00, 0,00 & 0,7 \\
Complexidade estatística & primária & 1,00, 1,00, 1,00 & 0,10, 0,10 & 0 \\
Entropia de permutação & primária & 1,00, 1,00, 1,00 & 0,10, 0,10 & 0 \\
Erro de previsão não linear & primária & 1,00, 1,00, 1,00 & 0,10, 0,00 & 0 \\
RQA · determinismo & exploratória & 1,00, 1,00, 1,00 & 0,10, 0,00 & 0,2 \\
RQA · laminaridade & exploratória & 1,00, 1,00, 1,00 & 0,00, 0,00 & 0 \\
Lyapunov (Wolf) & exploratória & 0,50, 1,00, 0,95 & 0,05, 0,00 & 0,2 \\
S-map (ganho não linear) & primária & 1,00, 1,00, 1,00 & 0,00, 0,00 & 0,2 \\
\bottomrule
\end{tabular}
\end{table}
